% !TEX root = ../main.tex
% =========================================================
%  CAPÍTULO 3 — MARCO TECNOLÓGICO
% =========================================================
 
\chapter{Capítulo tercero - Marco tecnológico}
 
La elección de las tecnologías que dan soporte a TankGo no ha sido arbitraria sino que responde
a un análisis previo de alternativas en el que se han ponderado criterios de
rendimiento, madurez del ecosistema, compatibilidad con el modelo de despliegue
cloud-native, coste operativo y alineación con las competencias adquiridas durante el
grado. Este capítulo documenta dicho análisis y justifica las decisiones adoptadas,
organizando las tecnologías según la capa arquitectónica en la que operan.
 
% ---------------------------------------------------------
\section{Capa de presentación: React y Leaflet.js}
% ---------------------------------------------------------
 
\subsection{React como framework de interfaz de usuario}
 
El cliente web de TankGo está desarrollado con \textbf{React 19}\cite{react}, la biblioteca de
interfaz de usuario de Meta. React adopta un paradigma declarativo basado en
componentes reutilizables y un DOM virtual que minimiza las operaciones de
renderizado sobre el árbol real, lo que resulta crítico en una aplicación como
TankGo, donde el mapa puede contener miles de marcadores que se actualizan
dinámicamente en función del viewport.
 
Entre las alternativas evaluadas se consideraron Vue.js y Angular \cite{vuejs,angular}. Angular, aunque
ofrece un entorno completo con \textit{dependency injection} integrada, introduce
una curva de aprendizaje significativa y un tamaño de bundle elevado para un
proyecto de alcance académico. Vue.js es una opción intermedia válida, pero su
ecosistema de integración con librerías cartográficas como Leaflet es menos maduro
que el de React, donde \texttt{react-leaflet}\cite{react_leaflet} ofrece una capa de abstracción directa sobre los componentes de ciclo de vida.
La prevalencia de React en el mercado laboral, alcanzando el 44,7\,\% según la encuesta anual de Stack Overflow de 2025, así 
como la familiaridad del autor con
la biblioteca, motivaron la elección final \cite{stackoverflow_survey_2025}.
 
\subsection{Leaflet.js como motor cartográfico}
 
La visualización geoespacial es clave en la experiencia de usuario de TankGo.
Para renderizar el mapa interactivo se evaluaron tres opciones principales: la API
de Google Maps, Mapbox GL JS y Leaflet.js.
 
La API de Google Maps queda descartada por su modelo de precios, que penaliza el
uso intensivo de marcadores y limita la personalización de las capas de datos sin
incurrir en costes adicionales \cite{logrocket_leaflet_guide}. Mapbox GL JS, si bien
ofrece renderizado WebGL de alto rendimiento y estilos vectoriales personalizables,
cambió su licencia de código abierto BSD a una licencia propietaria a partir de la
versión 2 en 2020, lo que excluye su uso libre en proyectos sin garantía de
continuidad de acceso \cite{geoapify_map_comparison}.
 
Leaflet.js es una biblioteca de código abierto bajo licencia BSD-2, pesa
aproximadamente 42\,KB en su versión minificada y ofrece un ecosistema de
\textit{plugins} especialmente extenso \cite{logrocket_leaflet_guide}. Su
integración nativa con teselas de OpenStreetMap elimina dependencias de
terceros propietarios, garantizando la continuidad del servicio sin costes de
API. El plugin \texttt{react-leaflet} \cite{react_leaflet} permite declarar capas y marcadores como
componentes React, integrándose de forma natural con el modelo de estado de la
aplicación. El soporte de agrupación de marcadores mediante \texttt{leaflet.markercluster}
resulta esencial para gestionar eficientemente las más de 11.000 estaciones de
servicio del catálogo.
 
% ---------------------------------------------------------
\section{Capa de gateway: Hono sobre Node.js}
% ---------------------------------------------------------
 
El API Gateway de TankGo está implementado con \textbf{Hono}, un framework
web ultraligero para entornos JavaScript modernos \cite{hono}. Hono se diseñó con compatibilidad
nativa para múltiples runtimes (Node.js, Deno, Bun, Cloudflare Workers) y ofrece
una API de enrutamiento expresiva junto con un rendimiento medido en benchmarks
independientes cercano a los 350.000 req/s sobre Node.js en escenarios de
\textit{hello world} \cite{hono_benchmarks}.
 
La alternativa natural habría sido Express.js, el framework más utilizado del
ecosistema Node.js\cite{express}. Sin embargo, Express es un framework sincrónico diseñado bajo
el modelo WSGI, mientras que Hono adopta la especificación Fetch API del estándar
web, lo que le permite gestionar peticiones concurrentes de forma asíncrona con
menor sobrecarga de memoria. Para un gateway cuya responsabilidad principal es
el proxy de peticiones HTTP (operación mayoritariamente I/O-bound),la naturaleza
asíncrona de Hono es especialmente adecuada.
 
% ---------------------------------------------------------
\section{Capa de servicios de negocio: FastAPI y Fastify}
% ---------------------------------------------------------
 
\subsection{FastAPI para servicios Python}
 
Los tres servicios de negocio implementados en Python (\texttt{gasolineras-service},
\texttt{recomendacion-service} y \texttt{prediction-service}) utilizan
\textbf{FastAPI} como framework HTTP. FastAPI está construido sobre
\textbf{Starlette} (capa ASGI) y \textbf{Pydantic v2} (validación de datos), lo que
le permite ofrecer soporte nativo de programación asíncrona sin imponer restricciones
al código de negocio síncrono.

Starlette es un microframework ASGI minimalista que proporciona el servidor ASGI, enrutamiento, middleware y 
utilidades para aplicaciones asíncronas \cite{starlette}. Pydantic permite definir modelos de datos mediante 
anotaciones de tipo en Python; valida y parsea automáticamente las entradas y salidas, generando errores legibles 
y esquemas JSON que FastAPI utiliza para documentar los modelos \cite{pydantic}. Gracias a esta combinación, 
FastAPI genera automáticamente un esquema OpenAPI por cada servicio, que se expone como documentación
interactiva (Swagger / Redoc) y facilita el consumo de la API por parte de clientes y herramientas de testing \cite{openapi}.
 
La elección de FastAPI frente a Flask y Django responde a tres razones principales.
En primer lugar, el rendimiento: los benchmarks disponibles sitúan a FastAPI ejecutado
sobre Uvicorn en el rango de 440 req/s en endpoints con lógica de negocio real, frente
a los aproximadamente 120 req/s de Flask en condiciones equivalentes, una diferencia
atribuible a la arquitectura ASGI frente a WSGI \cite{afolabi_fastapi_benchmark}. En
segundo lugar, la generación automática de documentación OpenAPI: FastAPI introspecciona
los \textit{type hints} de Python para generar el esquema OpenAPI sin código
adicional, lo que resulta crítico dado que el gateway agrega las especificaciones de
todos los servicios en un único portal Swagger unificado. En tercer lugar, la
integración con Pydantic garantiza la validación estricta de los parámetros de entrada
antes de que alcancen la lógica de negocio, reduciendo la superficie de ataque de la
API \cite{ijsrem_fastapi_comparison}.
 
Django\cite{django} habría añadido capas de abstracción innecesarias (sistema de plantillas,
panel de administración) que no son requeridas en servicios sin estado diseñados para
ser consumidos exclusivamente vía API REST. Flask\cite{flask}, aunque más ligero, carece de soporte
asíncrono nativo estable hasta su versión 3.0 y no genera documentación OpenAPI de
forma automática \cite{afolabi_fastapi_benchmark}.
 
\subsection{Fastify para el servicio de usuarios}
 
El servicio de autenticación y gestión de usuarios está implementado con
\textbf{Fastify}, el framework HTTP de alto rendimiento para Node.js. Fastify
es especialmente adecuado para este servicio por dos motivos: su sistema de plugins
permite integrar \texttt{@fastify/jwt} y \texttt{@fastify/cookie} con configuración
mínima, y su esquema de validación de peticiones basado en JSON Schema rechaza
\textit{payloads} malformados antes de que alcancen el controlador, sin necesidad de
bibliotecas adicionales.
 
La alternativa habría sido reimplementar el servicio de usuarios también en Python con
FastAPI, manteniendo homogeneidad tecnológica. Sin embargo, las operaciones de
autenticación (hashing con bcrypt, firma y verificación de JWT, gestión de cookies
httpOnly) están muy bien soportadas en el ecosistema Node.js con bibliotecas
como \texttt{bcryptjs} y \texttt{@fastify/jwt}, cuya madurez y documentación están
consolidadas. La decisión de mantener Node.js en este servicio refleja también el
criterio de utilizar la tecnología más idónea por capa, en lugar de imponer una
homogeneidad tecnológica artificial que habría penalizado el desarrollo.
 
% ---------------------------------------------------------
\section{Capa de datos: PostgreSQL y Apache Parquet}
% ---------------------------------------------------------
 
\subsection{PostgreSQL como base de datos relacional}
 
Todos los servicios que requieren persistencia utilizan \textbf{PostgreSQL}, el sistema
de gestión de bases de datos relacional de código abierto más avanzado de su categoría.
PostgreSQL ofrece soporte completo de transacciones ACID, tipos de datos geoespaciales
a través de la extensión \textbf{PostGIS} (utilizada en el servicio de recomendación
para operaciones como \texttt{ST\_DWithin}, \texttt{ST\_LineLocatePoint} y
\texttt{ST\_Distance}) y un modelo de concurrencia MVCC que minimiza los bloqueos en
lecturas simultáneas.
 
 La base de datos se aloja en \textbf{Neon}\cite{neon_tech2026}, una plataforma de PostgreSQL serverless en
 la nube que escala automáticamente a cero cuando no hay actividad, lo que resulta
idóneo para un proyecto académico con patrones de uso variables. Neon expone la misma
interfaz de conexión estándar de PostgreSQL, lo que permite utilizar los drivers
nativos \texttt{psycopg2} (Python) y \texttt{pg} (Node.js) sin adaptaciones.
 
\subsection{Apache Parquet para el pipeline de machine learning}
 
Los datos históricos de precios de combustibles exportados para el entrenamiento de
modelos se almacenan en formato \textbf{Apache Parquet}. Parquet es un formato de
almacenamiento columnar de código abierto, desarrollado originalmente por Twitter y
Cloudera, diseñado específicamente para cargas de trabajo analíticas
\cite{databricks_parquet}.
 
A diferencia de los formatos orientados a filas como CSV o JSON, Parquet agrupa los
valores de una misma columna de forma contigua en disco. Esta organización permite
a los motores analíticos leer únicamente las columnas necesarias para cada consulta,
reduciendo sustancialmente el volumen de I/O y mejorando la eficiencia de la
compresión, ya que los valores homogéneos dentro de una columna comprimen mejor que
filas heterogéneas \cite{datacamp_parquet}. En la práctica, los ficheros Parquet
pueden ser entre un 80 y un 90\,\% más pequeños que sus equivalentes en CSV para
conjuntos de datos tabulares de tipo analítico \cite{bckinfo_parquet}.
 
Para el pipeline de TankGo, en el que el modelo LightGBM consume únicamente
determinadas columnas del histórico (tipo de combustible, precio, fecha, coordenadas
y provincia), el formato columnar elimina la necesidad de deserializar filas completas,
acelerando significativamente la fase de carga de datos de entrenamiento. Los ficheros
Parquet se almacenan en \textbf{Google Cloud Storage} y son accesibles desde el
\texttt{prediction-service} mediante la biblioteca \texttt{google-cloud-storage}
de Python y \texttt{PyArrow} como motor de lectura.
 
% ---------------------------------------------------------
\section{Capa de inteligencia: LightGBM}
% ---------------------------------------------------------
 
%El módulo de predicción de precios de combustible utiliza \textbf{LightGBM},
%una implementación de \textit{Gradient Boosting} sobre árboles de decisión desarrollada
%por Microsoft. LightGBM introduce dos optimizaciones fundamentales respecto a
%implementaciones clásicas como XGBoost: \textit{Gradient-based One-Side Sampling}
%(GOSS), que descarta muestras con bajo gradiente para acelerar el entrenamiento, y
%\textit{Exclusive Feature Bundling} (EFB), que agrupa características mutuamente
%exclusivas para reducir la dimensionalidad efectiva \cite{tandfonline_lightgbm}.
 
%La elección de LightGBM para la predicción de series temporales de precios está
%respaldada por literatura académica reciente. Hartanto et al. demostraron que LightGBM
%supera o iguala a XGBoost en la predicción de precios de activos financieros con
%series temporales, obteniendo coeficientes $R^2$ superiores a 0,92 en conjuntos de
%prueba \cite{hartanto2023lightgbm}. Los precios de los combustibles, aunque no son
%exactamente equivalentes a los precios bursátiles, comparten la naturaleza de series
%temporales con dependencias estacionales y correlación con variables externas —en el
%caso de TankGo, el precio del barril de Brent obtenido mediante \texttt{yfinance}—,
%lo que hace aplicables las mismas consideraciones metodológicas.
 
%LightGBM se utiliza en modo de \textbf{regresión cuantílica}, entrenando tres modelos
%por estación: percentil 10 (límite inferior), mediana (predicción central) y
%percentil 90 (límite superior). Este enfoque permite expresar la incertidumbre de la
%predicción de forma calibrada, sin asumir distribuciones normales sobre los errores.
%Frente a modelos de series temporales clásicos como ARIMA o SARIMA, LightGBM admite
%directamente variables exógenas adicionales (día de la semana, mes, provincia, tipo de
%combustible) sin necesidad de transformaciones explícitas, lo que simplifica el
%pipeline de entrenamiento \cite{joiv_lightgbm}.
 
% ---------------------------------------------------------
\section{Capa de enrutamiento geoespacial: OpenRouteService}
% ---------------------------------------------------------
 
El servicio de recomendación de gasolineras en ruta requiere calcular geometrías de
trayecto y matrices de distancias reales entre puntos. Para ello se integra
\textbf{OpenRouteService} (ORS), una plataforma de enrutamiento de código abierto
desarrollada originalmente en el departamento GIScience de la Universidad de
Heidelberg, que ofrece una API REST con soporte para múltiples perfiles de vehículo
y la \textit{Matrix API} necesaria para el cálculo de desvíos reales desde la ruta
principal hasta cada candidato de repostaje \cite{gisops_routing_overview}.
 
La alternativa más directa era \textbf{OSRM} (\textit{Open Source Routing Machine}),
un motor de enrutamiento implementado en C++ con tiempos de respuesta en el rango de
los milisegundos para rutas regionales \cite{osm_wiki_routers}. Sin embargo, OSRM
requiere un servidor propio con requisitos de RAM elevados para el grafo de la red
viaria española, lo que es inviable en un entorno cloud serverless como Cloud Run,
donde los contenedores no mantienen estado persistente entre peticiones. ORS, en
cambio, expone una API REST gestionada que no requiere infraestructura adicional,
ofreciendo las funcionalidades de \textit{directions} y \textit{matrix} necesarias
para el algoritmo de recomendación \cite{gisops_routing_overview}.
 
Adicionalmente, la API de ORS soporta la evitación de
peajes como parámetro de la petición (incorporado en TankGo como opción
\texttt{evitar\_peajes}), funcionalidad que OSRM no contempla en su interfaz estándar.
 
%La integración técnica en \texttt{recomendacion-service} sigue el siguiente flujo:
%(1) ORS proporciona la geometría completa de la ruta mediante la API de
%\textit{directions}; (2) Shapely construye un corredor poligonal alrededor de esa
%geometría filtrando las estaciones candidatas dentro del radio de desvío configurado;
%(3) la \textit{Matrix API} de ORS calcula los desvíos reales desde la ruta hasta
%cada candidato; (4) un algoritmo de puntuación ponderada combina precio y desvío para
%obtener el ranking final de recomendaciones.
 
% ---------------------------------------------------------
\section{Capa de cartografía de datos: OpenStreetMap y Nominatim}
% ---------------------------------------------------------
 
Todos los datos cartográficos de base de TankGo (teselas del mapa, geocodificación
de búsquedas y geocodificación inversa) provienen de \textbf{OpenStreetMap} (OSM),
el proyecto colaborativo de cartografía libre que constituye la mayor base de datos
geoespacial de código abierto del mundo. Las teselas se consumen directamente desde
los servidores públicos de OSM a través de Leaflet.js.
 
La geocodificación (conversión de nombres de lugar en coordenadas y viceversa) se
realiza mediante \textbf{Nominatim}, el motor de búsqueda geográfica oficial de OSM.
El API Gateway actúa como proxy de las peticiones del frontend hacia Nominatim,
evitando que el navegador realice llamadas directas que podrían ser bloqueadas por la
política CORS de los servidores de OSM y centralizando el control de la caché de
respuestas.
 
La alternativa comercial habría sido la API de Geocoding de Google, cuyo modelo de
precios penaliza el número de peticiones a partir del umbral gratuito, introduciendo
una dependencia de proveedor incompatible con los objetivos de independencia
tecnológica del proyecto. La Geocoding API de Mapbox ofrece un modelo de precios más
accesible pero introduce igualmente dependencia de un servicio propietario.
 
% ---------------------------------------------------------
\section{Capa de inteligencia conversacional: Gemini y Google GenAI}
% ---------------------------------------------------------
 
El asistente de voz de TankGo implementa un pipeline
\textbf{STT $\rightarrow$ LLM con \textit{tool calling} $\rightarrow$ TTS}
íntegramente sobre \textbf{Gemini 2.5 Flash} de Google, accedido mediante la
biblioteca \texttt{@google/genai} y llamadas HTTP REST al endpoint
\texttt{/v1/models/gemini-2.5-flash:generateContent}.
 
La elección de un único modelo multimodal para las tres fases del pipeline —en lugar
de combinar, por ejemplo, Whisper (STT) + GPT-4o (LLM) + un motor TTS independiente—
simplifica significativamente la arquitectura del servicio, reduce la latencia
acumulada entre etapas y elimina la necesidad de gestionar múltiples claves de API
y contratos de servicio. Gemini 2.5 Flash está optimizado para latencia baja en
tareas de diálogo, soporta audio de entrada en base64 directamente en el payload de
 la petición, y su capacidad de \textit{tool calling} permite declarar herramientas
 con esquema JSON que el modelo invoca autónomamente para consultar precios y
 estaciones cercanas a través del gateway \cite{Gartner2026Trends}.
 
El \textit{tool calling} del asistente expone dos herramientas al modelo:
\texttt{get\_prices}, que consulta precios actuales de combustible por zona, y
\texttt{get\_nearest\_stations}, que devuelve las estaciones más cercanas a
unas coordenadas dadas. Ambas herramientas delegan en los endpoints HTTP del
gateway, de modo que el asistente actúa como interfaz conversacional del mismo
sistema de datos que consume el frontend web.
 
% ---------------------------------------------------------
\section{Plataforma de despliegue: Google Cloud}
% ---------------------------------------------------------
 
El despliegue en producción de TankGo utiliza tres servicios de \textbf{Google Cloud
Platform} (GCP): \textbf{Cloud Run}, \textbf{Cloud Build} y \textbf{Cloud Storage}\cite{cloud_run2026,cloud_build2026,cloud_storage2026}.
 
\textbf{Cloud Run} es una plataforma de cómputo serverless que ejecuta contenedores
Docker sin requerir gestión de infraestructura. Su modelo de escalado automático
—incluyendo la posibilidad de escalar a cero instancias cuando no hay tráfico—
 elimina los costes fijos de servidores permanentes, lo que es especialmente adecuado
 para un proyecto académico con patrones de uso irregulares \cite{cloud_run2026,Gartner2026Trends}.
Cada microservicio de TankGo se despliega como un servicio Cloud Run independiente
en la región \texttt{europe-west1}, garantizando aislamiento de fallos y escalado
independiente por servicio.
 
\textbf{Cloud Build} gestiona el pipeline de CI/CD mediante siete ficheros de
configuración independientes (\texttt{cloudbuild-*.yaml}), uno por servicio. Cada
pipeline construye la imagen Docker, la publica en \textbf{Artifact Registry} y
despliega la nueva revisión en Cloud Run. Los secretos de producción —claves de API,
credenciales de base de datos— se inyectan desde \textbf{Secret Manager} en tiempo
de despliegue, nunca almacenándose en el repositorio de código.
 
\textbf{Cloud Storage} actúa como almacén de objetos para los artefactos del pipeline
de machine learning: los ficheros Parquet con el histórico de precios exportados
diariamente por \texttt{gasolineras-service} y los modelos LightGBM serializados
generados por \texttt{prediction-service}.
 
La alternativa más directa habría sido \textbf{AWS} con Lambda o ECS. Sin embargo,
GCP ofrece una integración nativa entre Cloud Run, Gemini y las APIs de Google (Maps,
Geocoding, OAuth) que simplifica la configuración de identidades y permisos mediante
IAM, reduciendo la complejidad operativa del proyecto. Adicionalmente, los créditos
académicos del programa Google for Startups permitieron cubrir los costes de
infraestructura durante el período de desarrollo.
 
% ---------------------------------------------------------
\section{Síntesis: justificación del stack tecnológico}
% ---------------------------------------------------------
 
La Tabla~\ref{tab:stack_tecnologico} resume las tecnologías empleadas, su capa de
aplicación y el criterio principal que motivó su elección frente a las alternativas
evaluadas.
 
\begin{table}[H]
\centering
\caption{Resumen del stack tecnológico de TankGo}
\label{tab:stack_tecnologico}
\small
\begin{tabular}{llll}
\toprule
\textbf{Tecnología} & \textbf{Capa} & \textbf{Alternativas evaluadas} & \textbf{Criterio de elección} \\
\midrule
React 19            & Frontend      & Vue.js, Angular              & Ecosistema Leaflet, adopción \\
Leaflet.js          & Frontend      & Google Maps, Mapbox          & Licencia abierta, sin coste \\
Hono 4.10           & Gateway       & Express.js, Fastify          & Rendimiento async, WS nativo \\
FastAPI 0.115       & Backend Py    & Flask, Django                & ASGI, OpenAPI automático \\
Fastify 5.7         & Backend Node  & Express.js, NestJS           & JWT/cookie ecosystem \\
PostgreSQL / Neon   & Datos         & MySQL, MongoDB               & PostGIS, serverless, ACID \\
Apache Parquet      & ML pipeline   & CSV, JSON, Avro              & Columnar, compresión, Python \\
LightGBM 4.3        & ML modelo     & XGBoost, ARIMA, LSTM         & Velocidad, quantile regression \\
ORS (OpenRouteService) & Ruteo     & OSRM, Google Directions      & API REST gestionada, Matrix  \\
OpenStreetMap       & Cartografía   & Google Maps, Mapbox          & Datos libres, sin coste API \\
Nominatim           & Geocodificación & Google Geocoding            & Gratuito, integrado con OSM \\
Gemini 2.5 Flash    & IA / Voz      & GPT-4o + Whisper, Gemini Pro & Multimodal, latencia baja \\
Google Cloud Run    & Despliegue    & AWS Lambda, Azure Container  & Escala a cero, integración GCP \\
\bottomrule
\end{tabular}
\end{table}
 
El conjunto de tecnologías elegidas comparte un denominador común: todas son de uso
libre, todas tienen documentación oficial exhaustiva y todas cuentan con comunidades
activas que garantizan su mantenimiento a largo plazo. Este criterio de
\textit{sostenibilidad tecnológica} ha primado deliberadamente sobre la mera
novedad, asegurando que la solución desarrollada pueda ser mantenida y extendida
con independencia de los ciclos comerciales de proveedores específicos.
 
